    \section{Проверка программного окружения на наборе EuRoC}\label{sec:euroc}

    Прежде чем исследовать накопление ошибки в условиях, характерных для БПЛА, необходимо убедиться, что выбранные системы собраны и настроены корректно, а конвейер подготовки данных и оценки траекторий работает правильно.
    Для этого все три системы запущены на наборе EuRoC~\cite{euroc}, для которого в литературе имеются опубликованные результаты, и полученные ошибки сопоставлены с ними.
    Эксперименты этой главы не моделируют потерю GNSS: их задача -- проверка работоспособности окружения.

    \subsection{Программное окружение}

    Каждая система собрана в отдельном контейнере (podman), что изолирует несовместимые зависимости (ROS~1 и ROS~2, разные версии OpenCV) и делает сборку воспроизводимой.
    Версии исходного кода зафиксированы (таблица~\ref{table:env}).

    \begin{table}[H]
        \caption{Состав программного окружения}
        \label{table:env}
        \small
        \centering
        \renewcommand{\arraystretch}{1.2}
        \begin{tabularx}{\textwidth}{|l|X|X|}
            \hline
            \textbf{Система} & \textbf{Версия} & \textbf{Базовый образ} \\
            \hline
            ORB-SLAM3 & коммит \texttt{4452a3c}, Pangolin v0.6 & Ubuntu 20.04, ROS Noetic \\
            \hline
            VINS-Fusion & коммит \texttt{be55a93} с адаптацией к OpenCV~4 & Ubuntu 20.04, ROS Noetic \\
            \hline
            OpenVINS & релиз v2.7 & Ubuntu 22.04, ROS~2 Humble \\
            \hline
            Оценка траекторий & \texttt{evo} 1.38.1 & Python 3.14 \\
            \hline
        \end{tabularx}
    \end{table}

    При сборке потребовались две доработки.
    VINS-Fusion написан под OpenCV~3, поэтому устаревшие константы C-интерфейса заменены на их аналоги из C++-интерфейса OpenCV~4.
    В примерах ORB-SLAM3 отключено окно визуализации, так как контейнер не имеет дисплея.
    Алгоритмическая часть систем не изменялась, использовались штатные конфигурации для набора EuRoC.
    Эксперименты выполнены на ноутбуке с процессором Intel Core i5-13500H без графического ускорителя.

    \subsection{Данные и порядок проведения эксперимента}

    Из набора EuRoC выбраны пять последовательностей разной сложности (таблица~\ref{table:euroc_seq}): три полёта в машинном зале с эталоном от лазерного трекера и два полёта в комнатах с системой захвата движения Vicon.
    Все системы работали в конфигурации «монокулярная камера и IMU» (левая камера \texttt{cam0}, 20~Гц; IMU 200~Гц), наиболее реалистичной для малого БПЛА.

    \begin{table}[H]
        \caption{Использованные последовательности EuRoC}
        \label{table:euroc_seq}
        \small
        \centering
        \renewcommand{\arraystretch}{1.2}
        \begin{tabularx}{\textwidth}{|l|c|c|c|X|}
            \hline
            \textbf{Последовательность} & \makecell{\textbf{Длитель-}\\\textbf{ность, с}} & \makecell{\textbf{Путь,}\\\textbf{м}} & \textbf{Кадров} & \textbf{Особенности} \\
            \hline
            MH\_01\_easy & 184 & 80 & 3682 & медленный полёт, хорошее освещение \\
            MH\_03\_medium & 135 & 130 & 2700 & быстрое движение \\
            MH\_05\_difficult & 114 & 97 & 2273 & быстрое движение, тёмные участки \\
            V1\_02\_medium & 86 & 76 & 1710 & резкие повороты в небольшой комнате \\
            V2\_02\_medium & 117 & 83 & 2348 & резкие повороты в небольшой комнате \\
            \hline
        \end{tabularx}
    \end{table}

    Подготовка данных и запуск автоматизированы скриптами:
    \begin{enumerate}
        \item последовательности загружаются в исходном формате (изображения и CSV-файлы) и для VINS-Fusion и OpenVINS преобразуются в формат rosbag для ROS~1 и ROS~2 с теми же названиями топиков, что в оригинальном наборе;
        \item каждая система запускается в своём контейнере, данные подаются \emph{в реальном времени}: если система не успевает обработать кадр до прихода следующего, кадр пропускается, как это происходило бы на борту;
        \item оценённая траектория корпуса (IMU) сохраняется в общем формате TUM (время, положение, кватернион);
        \item для VINS-Fusion сохраняются две траектории -- выход VIO и траектория после замыкания петель, -- чтобы оценить вклад замыкания петель отдельно.
    \end{enumerate}

    Оценка выполнена по методике раздела~\ref{sec:review}.
    Оценённые позы сопоставляются с эталоном по времени (допуск $10$~мс), траектория совмещается с эталоном преобразованием $SE(3)$ по методу Умэямы~\cite{umeyama}.
    Масштаб при совмещении не подбирается: при использовании IMU он наблюдаем, и его ошибка должна входить в метрику.
    Вычисляются ATE RMSE~(\ref{eq:ate}), максимальная абсолютная ошибка и относительный дрейф -- средняя ошибка перемещения на всех отрезках длиной $10$~м в процентах от длины отрезка.
    Такая методика совпадает с использованной в работе~\cite{orbslam3}, что позволяет сравнить результаты.

    \subsection{Результаты}

    Траектории всех систем показаны на рис.~\ref{fig:euroc_traj}, значения ATE и их сравнение с опубликованными результатами приведены в таблице~\ref{table:euroc_ate}, относительный дрейф -- в таблице~\ref{table:euroc_drift}.
    Опубликованные значения взяты из сравнительной таблицы работы~\cite{orbslam3}, где для ORB-SLAM3 приведена медиана по десяти запускам, а VINS-Mono запущен авторами со штатной конфигурацией; для OpenVINS опубликованных результатов в этой конфигурации не найдено.

    \begin{figure}[H]
        \centering
        \includegraphics[width=\textwidth]{euroc_traj.png}
        \caption{Траектории на последовательностях EuRoC после совмещения с эталоном (вид сверху)}
        \label{fig:euroc_traj}
    \end{figure}

    \begin{table}[H]
        \caption{ATE RMSE на последовательностях EuRoC, см (монокулярная камера и IMU)}
        \label{table:euroc_ate}
        \small
        \centering
        \renewcommand{\arraystretch}{1.2}
        \begin{tabular}{|l|c|c|c|c|c|}
            \hline
            \textbf{Метод} & \textbf{MH01} & \textbf{MH03} & \textbf{MH05} & \textbf{V102} & \textbf{V202} \\
            \hline
            ORB-SLAM3 & 2,4 & 4,7 & 7,1 & 1,9 & 2,7 \\
            ORB-SLAM3, опубликовано~\cite{orbslam3} & 6,2 & 4,6 & 5,7 & 1,5 & 2,1 \\
            \hline
            VINS-Fusion (VIO) & 18,5 & 12,7 & 32,6 & 17,5 & 11,4 \\
            VINS-Fusion с замыканием петель & 6,7 & 6,3 & 15,0 & 10,6 & 10,0 \\
            VINS-Mono, по данным~\cite{orbslam3} & 8,4 & 7,4 & 14,7 & 6,6 & 9,0 \\
            \hline
            OpenVINS & 10,4 & 11,1 & 44,0 & 6,7 & 6,7 \\
            \hline
        \end{tabular}
    \end{table}

    \begin{table}[H]
        \caption{Относительный дрейф на отрезках 10~м, \%}
        \label{table:euroc_drift}
        \small
        \centering
        \renewcommand{\arraystretch}{1.2}
        \begin{tabular}{|l|c|c|c|c|c|}
            \hline
            \textbf{Метод} & \textbf{MH01} & \textbf{MH03} & \textbf{MH05} & \textbf{V102} & \textbf{V202} \\
            \hline
            ORB-SLAM3 & 0,89 & 1,34 & 1,09 & 0,96 & 0,41 \\
            VINS-Fusion (VIO) & 1,94 & 1,80 & 3,87 & 1,81 & 1,32 \\
            VINS-Fusion с замыканием петель & 1,54 & 1,57 & 3,48 & 1,41 & 1,19 \\
            OpenVINS & 1,06 & 1,63 & 2,81 & 1,09 & 0,69 \\
            \hline
        \end{tabular}
    \end{table}

    \subsection{Анализ результатов}

    \textbf{Соответствие опубликованным результатам.}
    Ошибки ORB-SLAM3 совпадают с опубликованными в пределах $1$--$2$~см на всех последовательностях, на MH01 полученная ошибка меньше опубликованной.
    Ошибки VINS-Fusion с замыканием петель близки к опубликованным для VINS-Mono на четырёх из пяти последовательностей; на V102 ошибка больше ($10{,}6$ против $6{,}6$~см).
    Расхождения такого масштаба ожидаемы: в работе~\cite{orbslam3} приведена медиана десяти запусков, а здесь -- один запуск, при этом многопоточные системы дают разные результаты от запуска к запуску; кроме того, VINS-Fusion в монокулярно-инерциальном режиме лишь повторяет VINS-Mono, а не совпадает с ним полностью.
    Таким образом, сборка и настройка систем, подготовка данных и расчёт метрик признаны корректными.

    \textbf{Сравнение систем.}
    Наименьшую ошибку на всех последовательностях показал ORB-SLAM3 (ATE $1{,}9$--$7{,}1$~см, дрейф $0{,}4$--$1{,}3$\%).
    Замыкание петель уменьшает ошибку VINS-Fusion в $1{,}1$--$2{,}8$ раза, причём сильнее всего (в $2$--$2{,}8$ раза) на полётах в машинном зале с повторными пролётами над одними и теми же участками.
    OpenVINS без замыкания петель точнее VIO-части VINS-Fusion на четырёх последовательностях из пяти, но уступает ORB-SLAM3.
    Наиболее трудной оказалась MH05: ошибка всех систем на ней максимальна, а у OpenVINS достигает $44$~см.
    Порядок систем согласуется с ожиданиями главы~\ref{sec:methods}: SLAM с замыканием петель точнее одометрии; при этом фильтр OpenVINS с долгоживущими SLAM-точками в большинстве случаев не уступает оптимизации в скользящем окне.

    \textbf{Различия в начале оценки.}
    Системы начинают выдавать оценку в разные моменты.
    Последовательности машинного зала начинаются с неподвижного аппарата, и OpenVINS со статической инициализацией начинает работу только после начала движения: на MH01 оценка охватывает $138$ из $184$~с.
    Максимальная ошибка ORB-SLAM3 ($17$--$58$~см) существенно превышает его среднюю ошибку; предположительно, она приходится на первые секунды до завершения инерциальной инициализации, когда масштаб ещё не определён.
    Для передачи оценки в автопилот это означает, что выход системы необходимо отбрасывать до завершения инициализации.

    \textbf{Сбой VINS-Fusion.}
    При первом запуске на V2\_02 оценка VINS-Fusion разошлась через $5$~с после инициализации: решатель нелинейных наименьших квадратов перестал находить шаг оптимизации, и ошибка положения за время полёта выросла до десятков километров.
    Два повторных запуска на той же последовательности завершились штатно (в таблицах приведён последний из них).
    Это указывает на случайный сбой инициализации, а не на ошибку настройки.
    Для навигации расходящаяся оценка опаснее неточной, поэтому в дальнейших экспериментах помимо средней ошибки необходимо фиксировать долю запусков, завершившихся расхождением.

    \subsection{Выводы по главе}

    Создано воспроизводимое программное окружение для запуска ORB-SLAM3, VINS-Fusion и OpenVINS и единый конвейер подготовки данных и оценки траекторий.
    Проверка на пяти последовательностях EuRoC показала, что все три системы работают в реальном времени в конфигурации «монокулярная камера и IMU», а полученные ошибки ORB-SLAM3 и VINS-Fusion соответствуют опубликованным результатам, что подтверждает корректность окружения.

    Проверка выявила свойства систем, важные для дальнейшего исследования: различное время инициализации, повышенную ошибку в её начале и возможность расхождения оценки.
    Ограничением проведённой проверки является один запуск на последовательность; в основной серии экспериментов каждая конфигурация будет запускаться многократно.
    Кроме того, траектории EuRoC короткие и получены в помещении, поэтому исследование накопления ошибки после потери GNSS в зависимости от высоты, скорости и типа поверхности требует данных, приближенных к условиям полёта БПЛА.
    \clearpage
